% Imported from LaTeX_and_pdfs/AGAT_dump.tex; see provenance.json.
\part{Reading notes on Cox, Little, O'Shea's \emph{Ideals, Varieties, and Algorithms: An Introduction to Computational Algebraic Geometry and Commutative Algebra}}


\section{Geometry, Algebra, and Algorithms}

\subsection{Polynomials and Affine Space}

fields are important is that linear algebra works over \emph{any} field

\begin{definition}[2] set of all polynomials in $x_1 , \dots , x_n$ with coefficients in $k$, denoted $k[x_1, \dots , x_n]$

\end{definition}

polynomial $f$ \emph{divides} polynomial $g$ provided $g= fh$ for some $h \in k[x_1, \dots , x_n ]$

$k[x_1, \dots, x_n]$ satisfies all field axioms except for existence of multiplicative inverses; commutative ring, $k[x_1, \dots , x_n]$ \emph{polynomial ring}


\subsubsection*{Exercises for 1}

\exercisehead{1}
$\mathbb{F}_2$ commutative ring since it's an abelian group under addition, commutative in multiplication, and multiplicative identity exists, namely $1$.  It is a field since for $1\neq 0$, the multiplicative identity is $1$.  

\exercisehead{2}
\begin{enumerate}
\item[(a)]
\item[(b)]
\item[(c)]
\end{enumerate}




\subsection{Affine Varieties }


\subsection{Parametrizations of Affine Varieties}


\subsection{Ideals}



\subsection{Polynomials of One Variable}




\section{Groebner Bases}

\subsection{Introduction}




\subsection{Orderings on the Monomials in $k[x_1, \dots , x_n]$ }




\subsection{A Division Algorithm in $k[x_1, \dots , x_n ]$ }



\subsection{Monomial Ideals and Dickson's Lemma }


\subsection{The Hilbert Basis Theorem and Groebner Bases}


\subsection{Properties of Groebner Bases}



\subsection{Buchberger's Algorithm}






\section{Elimination Theory}



\subsection{The Elimination and Extension Theorems}


\subsection{The Geometry of Elimination}



\section{The Algebra-Geometry Dictionary}


\subsection{Hilbert's Nullstellensatz}


\subsection{Radical Ideals and the Ideal-Variety Correspondence}



\section{Polynomial and Rational Functions on a Variety}


\subsection{Polynomial Mappings }


\section{Robotics and Automatic Geometric Theorem Proving}



\subsection{Geometric Description of Robots}






